%% algorithm_description.tex
%%
%% A high-level description of the algorithms in MinimalConfusableSets,
%% reconstructed primarily from reading the source code.
%% Written with Claude's assistance, May 2026.
%%
%% WARNING: This document was reconstructed from code, not from the author's
%% own notes. Statements marked with \textbf{[?]} indicate places where
%% the intent has been inferred rather than read directly.

\documentclass[12pt,a4paper]{article}

\usepackage{amsmath,amssymb,amsthm}
\usepackage{hyperref}
\usepackage{geometry}
\geometry{margin=2.5cm}

\newtheorem{definition}{Definition}
\newtheorem{remark}{Remark}
\newtheorem{observation}{Observation}

\title{Minimal Confusable Multisets:\\
A Description of the Algorithms in \texttt{MinimalConfusableSets}}

\author{Lester (code), Claude (document reconstruction)}
\date{May 2026}

\begin{document}
\maketitle

\begin{abstract}
This document gives a high-level description of the algorithms implemented in
the \texttt{MinimalConfusableSets} subdirectory of the Echinocoder repository.
The goal of that code is to search for, and determine the minimum size of,
\emph{confusable multisets}: pairs of equal-size multisets of vectors in
$\mathbb{R}^k$ that cannot be distinguished by a given family of encoding
directions. The description here has been reconstructed from reading the source
code rather than from the author's own notes; uncertain inferences are flagged
with \textbf{[?]}.
\end{abstract}

\tableofcontents
\newpage

%% -----------------------------------------------------------------------
\section{Background and motivation}
\label{sec:background}

\subsection*{The Echinocoder project}

The broader Echinocoder project is concerned with \emph{embeddings} of
multisets of vectors: injective continuous maps from the space
$SP^n(\mathbb{R}^k)$ of $n$-element multisets of vectors in $\mathbb{R}^k$
into some Euclidean space $\mathbb{R}^N$. A fundamental requirement is that
such a map be invariant under all permutations of the $n$ input vectors whilst
still distinguishing genuinely distinct multisets.

\subsection*{Dotting encoders}

One natural class of encoders---the \emph{dotting encoders}---works as follows.
Fix $m$ direction vectors $\vec{d}_1, \ldots, \vec{d}_m \in \mathbb{R}^k$,
called the \emph{encoding directions} (or \emph{good bats} in the project's
terminology). For an input multiset $X = \{\vec{v}_1, \ldots, \vec{v}_n\}
\subset \mathbb{R}^k$, and for each encoding direction $\vec{d}_j$ separately,
compute the $n$ dot products $\vec{v}_1\cdot\vec{d}_j,\ldots,
\vec{v}_n\cdot\vec{d}_j$ and sort them into non-decreasing order.
Concatenating these $m$ sorted $n$-tuples gives the encoding of $X$ as a
vector in $\mathbb{R}^{mn}$.  This output is automatically symmetric (invariant
under permutations of the $\vec{v}_i$).  The encoder is \emph{injective}---and
hence an embedding---if and only if no two distinct multisets in
$SP^n(\mathbb{R}^k)$ produce the same output.

The central open question for dotting encoders is: how many directions $m$ are
needed to guarantee injectivity for multisets of size $n$ in dimension $k$, and
can $m$ be kept small?

\subsection*{Adversarial inputs: bad bats and confusable multisets}

To probe the limits of any dotting encoder one can try to construct pairs of
\emph{distinct} multisets that it cannot distinguish.  The
\texttt{MinimalConfusableSets} sub-project is devoted to this adversarial
construction, and in particular to finding the \emph{smallest} such pairs for
given $m$ and $k$.

The construction exploits a linear-algebraic obstruction that exists for
\emph{any} choice of encoding directions. In $\mathbb{R}^k$, any $k-1$
linearly independent vectors span a $(k-1)$-dimensional hyperplane, and there
is exactly one direction in $\mathbb{R}^k$ (up to scaling) orthogonal to all
of them. Given $m$ encoding directions, partition them into
$M = \lceil m/(k-1) \rceil$ groups of (at most) $k-1$ directions each. For
the $j$-th group, take the unique direction in $\mathbb{R}^k$ orthogonal to
all members of that group; call it $\vec{q}_j$. These $M$ vectors
$\vec{q}_1, \ldots, \vec{q}_M \in \mathbb{R}^k$ are the \emph{bad bats}.

Crucially, each bad bat $\vec{q}_j$ satisfies $\vec{q}_j \cdot \vec{d} = 0$
for the $k-1$ encoding directions $\vec{d}$ in its own group, but is
generically \emph{not} orthogonal to the encoding directions in any other
group. In particular, once $m \geq k$ directions are in general position, no
non-zero vector in $\mathbb{R}^k$ can be orthogonal to \emph{all} of them
simultaneously: bad bats exist because each is orthogonal only to a small
group of $k-1$ directions, not to the full set of $m$.

\subsection*{Why the even-odd construction yields confusable pairs}

Building multisets from scaled sums of the bad bats and applying the
even-odd hypercube construction (Section~\ref{sec:even-odd}) yields pairs of
multisets $(EE, OO)$ that are \emph{confusable}: for \emph{every} encoding
direction $\vec{d}_j$, the multiset of dot products
$\{\,\vec{v} \cdot \vec{d}_j : \vec{v} \in EE\,\}$ equals
$\{\,\vec{v} \cdot \vec{d}_j : \vec{v} \in OO\,\}$.

The key reason this works is a \emph{toggling symmetry}. Consider any
encoding direction $\vec{d}$ belonging to group $g$, so that
$\vec{q}_g \cdot \vec{d} = 0$. For any subset sum
$\vec{s} = \sum_{j \in I} \alpha_j \vec{q}_j$ of the bad bats, the dot
product $\vec{s} \cdot \vec{d} = \sum_{j \in I} \alpha_j\,(\vec{q}_j \cdot
\vec{d})$ receives no contribution from the $j = g$ term.  Consequently,
including or excluding $\vec{q}_g$ from the sum leaves $\vec{s} \cdot \vec{d}$
unchanged---while flipping the parity of $|I|$, which is what distinguishes
the even ($E$) and odd ($O$) families.  The dot-product values of $E$ and $O$
with respect to $\vec{d}$ are therefore identical as multisets, and this holds
simultaneously for every encoding direction (each of which belongs to some
group and has the corresponding bad bat orthogonal to it).

\subsection*{The goal of \texttt{MinimalConfusableSets}}

The size $|EE| = |OO|$ of a confusable pair depends on the specific bad-bat
configuration and on which subset-sum coincidences are imposed.  The
\texttt{MinimalConfusableSets} code searches for the \emph{smallest} confusable
pairs achievable for given values of $M$ and $k$.  A confusable pair of size
$s$ demonstrates that no dotting encoder with approximately $M(k-1)$ encoding
directions can injectively encode all multisets of size $\geq s$, regardless
of which specific directions are chosen.

The two primary parameters throughout this document are therefore:
\begin{itemize}
  \item $k$: the dimension of the ambient space,
  \item $M$: the number of bad-bat vectors (equivalently, the number of
    groups of encoding directions).
\end{itemize}

\subsection*{An important caveat: the construction is not known to be exhaustive}

The even-odd hypercube construction gives a \emph{sufficient} condition for
producing confusable pairs, but it is not known to be \emph{necessary}.  The
definition of a confusable pair is simply: two distinct multisets
$X, Y \subset \mathbb{R}^k$ of equal size such that for every encoding
direction $\vec{d}_j$ the multiset of dot products
$\{\,\vec{v}\cdot\vec{d}_j : \vec{v}\in X\,\}$ equals
$\{\,\vec{v}\cdot\vec{d}_j : \vec{v}\in Y\,\}$.
Nothing in this definition requires $X$ or $Y$ to arise from subset sums of
bad bats, or for their elements to have any particular combinatorial structure.
Sporadic confusable pairs may exist---pairs whose elements bear no relationship
to any even-odd hypercube construction---and some of those might be smaller
than any pair the code can find.

Consequently the output of \texttt{MinimalConfusableSets} is an upper bound on
the true minimum confusable-pair size for two \emph{distinct} reasons:
\begin{enumerate}
  \item \emph{Incomplete search within the hypercube family.}  Some pruning
    strategies in the code are acknowledged to be potentially over-aggressive
    (Section~\ref{sec:pruning}), so not every hypercube-family pair is
    guaranteed to be found.  This is a correctable deficiency.
  \item \emph{The hypercube family may not cover all confusable pairs.}  Even
    a hypothetically perfect, bug-free search within the current framework
    would only establish a minimum over pairs of a specific structured form.
    The truly minimal confusable pair might take a completely different form
    that the construction never generates.  No theorem, and at present no
    strong conjecture, rules this out.
\end{enumerate}
Reason~(1) affects the reliability of the search within its chosen scope.
Reason~(2) is a more fundamental limitation on the scope itself.

%% -----------------------------------------------------------------------
\section{The even-odd hypercube construction}
\label{sec:even-odd}

Given $M$ vectors $\vec{q}_1, \ldots, \vec{q}_M \in \mathbb{R}^k$, define the
$2^M$ \emph{subset sums}
\[
  S_I \;=\; \sum_{i \in I} \vec{q}_i, \qquad I \subseteq \{1, \ldots, M\},
\]
and partition them according to the parity of $|I|$:
\begin{align*}
  E &= \{\, S_I : |I| \text{ is even}\,\},\\
  O &= \{\, S_I : |I| \text{ is odd}\,\}.
\end{align*}
Both $E$ and $O$ are multisets (a given point may appear with multiplicity $>1$
if different index sets $I$ produce the same sum). By a simple symmetry argument
$|E| = |O| = 2^{M-1}$ counting multiplicity.

Let $C$ be the ``collision'' multiset: for each point $\vec{v}$, its
multiplicity in $C$ is $\min(\text{mult}_{E}(\vec{v}),\,
\text{mult}_{O}(\vec{v}))$. Define the \emph{exclusive} parts
\[
  EE = E \setminus C, \qquad OO = O \setminus C,
\]
where subtraction is taken in the multiset sense. By construction
$|EE| = |OO|$.

\begin{definition}
  A \emph{confusable multiset pair} is a pair $(EE, OO)$ produced by the
  even-odd hypercube construction with $|EE| = |OO| > 0$.
\end{definition}

\begin{remark}
  The pair $(EE, OO)$ is confusable in the sense that both multisets have the
  same multiset of dot products with any direction orthogonal to all
  $\vec{q}_i$. In particular, any encoding that uses only such orthogonal
  directions cannot distinguish $EE$ from $OO$.
\end{remark}

\subsection*{Why $|EE|$ can be zero}

If $E$ and $O$ coincide exactly as multisets (total cancellation), then $C = E
= O$ and $EE = OO = \emptyset$. The code calls this outcome \emph{collapse},
and the corresponding configuration of bad bats is uninteresting. The search
is for configurations that \emph{avoid} collapse.

%% -----------------------------------------------------------------------
\section{Scaling the bad bats: unscaled and scaled representations}
\label{sec:scaling}

In the code, the bad-bat vectors are separated into \emph{directions}
$\vec{w}_1, \ldots, \vec{w}_M \in \mathbb{R}^k$ (fixed integer vectors in
``general position'') and \emph{scalings} $\alpha_1, \ldots, \alpha_M \in
\mathbb{R}$, so that
\[
  \vec{q}_i = \alpha_i \,\vec{w}_i.
\]
The \emph{unscaled bat matrix} $W$ is the $M \times k$ matrix whose $i$-th row
is $\vec{w}_i$. The \emph{scaled bat matrix} $B = \operatorname{diag}(\alpha)
\, W$ has $i$-th row $\alpha_i \vec{w}_i$.

The purpose of this split is to allow the problem of finding useful
configurations to be formulated as a linear algebra question (Section
\ref{sec:alpha-attacking}).

The unscaled directions are chosen randomly in general position, meaning every
subset of at most $k$ of the $\vec{w}_i$ is linearly independent. This
genericity ensures that the null-space calculation (Section
\ref{sec:alpha-attacking}) captures the right constraints.

%% -----------------------------------------------------------------------
\section{Vertex matches}
\label{sec:vertex-matches}

\begin{definition}
  A \emph{vertex match} is a tuple $\ell = (\ell_1, \ldots, \ell_M) \in
  \{-1, 0, +1\}^M$ such that the number of $+1$ entries is even and the number
  of $-1$ entries is odd.
\end{definition}

Geometrically, a vertex match encodes a coincidence between two subset sums:
the $+1$ positions form one summand and the $-1$ positions form the other, so
that the constraint $\sum_i \ell_i \vec{q}_i = 0$ expresses the fact that
those two subsets produce the same point. The parity condition (even $+1$s, odd
$-1$s) ensures that one summand contributes to $E$ and the other to $O$,
producing the kind of coincidence needed for a non-trivial confusable pair.
\textbf{[?]}

\begin{definition}
  A vertex match $\ell$ is \emph{useful} (for given $k$) if it has at least
  $k+1$ non-zero entries. The reason is that the sum of $\leq k$ linearly
  independent vectors in $k$-dimensional space is always non-zero, so a match
  with only $k$ or fewer active positions cannot hold for generic bad bats.
\end{definition}

\subsection*{Tristate vs bistate vertex matches}

Two variants of vertex match are present in the code:
\begin{itemize}
  \item \textbf{Tristate}: entries in $\{-1, 0, +1\}$ with the parity
    constraint above. This is the original and more general formulation.
  \item \textbf{Bistate}: entries in $\{0, +1\}$ only (no $-1$ entries). This
    is a more recent variant \textbf{[?]} corresponding to the
    \texttt{NO\_MINUS\_ONES\_BRANCH} in git history. The geometric meaning of
    the bistate constraint is that a single subset sums to zero, $\sum_{i \in
    I} \vec{q}_i = 0$, rather than two subsets having equal sums.
\end{itemize}
At runtime exactly one variant is active, selected via \texttt{config.py}.

\subsection*{Equivalent places}

When the current partial L-matrix (see Section \ref{sec:L-matrices}) has not
yet broken any symmetry between some subset of the $M$ bad-bat positions, those
positions are \emph{equivalent}: permuting them within the next row cannot
produce anything inequivalent to what has already been generated. The code
tracks these equivalence classes explicitly via the \texttt{Equivalent\_Places}
class, and uses them to avoid redundant generation of permuted rows.

%% -----------------------------------------------------------------------
\section{L-matrices}
\label{sec:L-matrices}

An \emph{L-matrix} is a matrix $L$ whose rows are vertex matches. Each row
imposes one coincidence constraint; the full matrix imposes all of them
simultaneously.

The ultimate output of the search is a ``good'' L-matrix: one for which (a)
there exists a scaling $\alpha$ consistent with all the constraints, and (b)
the resulting confusable pair $(EE, OO)$ is non-empty and as small as possible.

%% -----------------------------------------------------------------------
\section{The alpha-attacking matrix and collapse detection}
\label{sec:alpha-attacking}

Substituting $\vec{q}_i = \alpha_i \vec{w}_i$ into the $r$-th vertex-match
constraint
\[
  \sum_{j=1}^M L_{rj}\, \vec{q}_j = \vec{0}
\]
and writing out the $k$ component equations gives, for each component $\kappa$:
\[
  \sum_{j=1}^M L_{rj} \,\alpha_j \, w_{j\kappa} \;=\; 0.
\]
Collecting all $R$ constraints and all $k$ components into a single linear
system $A \alpha = 0$ defines the \emph{alpha-attacking matrix}
\[
  A_{(r-1)k+\kappa,\; j} \;=\; L_{rj} \cdot w_{j\kappa},
  \qquad r=1,\ldots,R,\; \kappa=1,\ldots,k,\; j=1,\ldots,M.
\]
This is computed by \texttt{alpha\_attacking\_matrix} in
\texttt{confusable\_multisets.py}.

The null space $\mathcal{N}(A)$ is the set of valid scalings $\alpha$. Three
cases arise:
\begin{enumerate}
  \item $\dim \mathcal{N}(A) = 0$: the only solution is $\alpha = 0$; all bad
    bats vanish. This is \emph{complete collapse}; the matrix is discarded.
  \item $\dim \mathcal{N}(A) = M$: $A$ is the zero matrix, so $\alpha$ is
    completely free. Any all-ones vector lies in the null space: no collapse.
  \item $0 < \dim \mathcal{N}(A) < M$: check whether every coordinate
    direction $e_j$ is non-zero in at least one basis vector of
    $\mathcal{N}(A)$. If so, a linear combination exists with all $\alpha_j
    \neq 0$ (proved in the code comments, attributed to a OneNote entry). If
    some coordinate $j$ is zero in every basis vector, then $\alpha_j = 0$ is
    forced: bat $j$ collapses.
\end{enumerate}

\subsection*{Finding a non-collapsing point in the null space}

When the null space is non-trivial and non-collapsing, the code needs a
specific point $\alpha^* \in \mathcal{N}(A)$ with all $\alpha^*_j \neq 0$.
This is done by \texttt{nonzero\_lin\_comb.combine\_many}: starting from the
first basis vector, it combines basis vectors one at a time, at each step
choosing an integer multiplier $\lambda$ (the smallest integer avoiding a
finite set of ``forbidden'' values) that keeps all previously-established
non-zero coordinates non-zero while activating the next coordinate.

%% -----------------------------------------------------------------------
\section{Computing the confusable pair}
\label{sec:computing}

Given the scaled bat matrix $B = \operatorname{diag}(\alpha^*) W$, the code
computes $(EE, OO)$ via a \emph{meet-in-the-middle} approach:

\begin{enumerate}
  \item Split the $M$ rows of $B$ into a left half ($M/2$ rows) and a right
    half.
  \item For each half, enumerate all $2^{M/2}$ subset sums, keeping track of
    the parity of the subset size. This gives four counters: $L_{\text{even}}$,
    $L_{\text{odd}}$, $R_{\text{even}}$, $R_{\text{odd}}$.
  \item Convolve:
    \begin{align*}
      E &= (L_{\text{even}} * R_{\text{even}}) + (L_{\text{odd}} * R_{\text{odd}}),\\
      O &= (L_{\text{even}} * R_{\text{odd}}) + (L_{\text{odd}} * R_{\text{even}}),
    \end{align*}
    where $*$ denotes convolution of the counter-valued functions on
    $\mathbb{R}^k$.
  \item Compute $C$, $EE$, $OO$ as in Section \ref{sec:even-odd}.
\end{enumerate}

This reduces the cost from $O(2^M)$ to $O(2^{M/2})$ at the expense of the
convolution step.

%% -----------------------------------------------------------------------
\section{The search algorithm}
\label{sec:search}

The outer loop, driven by \texttt{generate\_viable\_vertex\_match\_matrices} in
\texttt{vertex\_matches.py}, enumerates L-matrices via a depth-first search:

\begin{enumerate}
  \item Start with the empty matrix.
  \item At each level of the DFS, extend the current matrix by one row, chosen
    from all valid vertex matches respecting the current equivalence classes of
    bat positions.
  \item Apply pruning tests (Section \ref{sec:pruning}) to decide whether to
    yield the current matrix and/or recurse deeper.
  \item For each yielded matrix, call the \emph{decider function} to test for
    non-collapse and, if non-collapsing, to compute $(EE, OO)$ and record
    $|EE|$.
  \item Track the smallest $|EE|$ found across all matrices and all choices of
    unscaled bat matrix.
\end{enumerate}

\subsection*{Robustness via multiple bat matrices}

A collapse result for a specific unscaled bat matrix $W$ does not rule out a
non-collapsing solution for a different $W$. To guard against this, the demo
driver (\texttt{tom\_demo.py}) runs multiple independent \texttt{Rational\_Decider}
instances, each using a differently-seeded random bat matrix. A matrix is
treated as collapsing only when all deciders agree. The comment in the code
contains a probabilistic argument (incomplete at the time of writing) that this
multi-decider approach can drive the false-collapse probability arbitrarily
low.

%% -----------------------------------------------------------------------
\section{Pruning strategies}
\label{sec:pruning}

Several pruning strategies are implemented to reduce the search space.

\subsection{Short-row pruning (\texttt{prune\_short\_rows})}

If the reduced row echelon form (RREF) of the current matrix contains a row
with between $1$ and $k$ non-zero entries, the matrix is discarded. The
reasoning is that such a row would force a linear combination of $\leq k$
specific bad-bat directions to vanish, which is impossible when those bats are
in general position. This pruning is the most effective and is on by default.

\subsection{Maximum-rows pruning (\texttt{prune\_max\_rows})}

The RREF of a viable L-matrix cannot have more than a certain number of rows,
determined by $M$ and $k$. The code can optionally stop deepening the DFS once
this maximum is reached. This pruning is marked as ``probably buggy/wrong'' in
the source.

\subsection{Pivot pruning (\texttt{prune\_pivots})}

Discards matrices whose RREF pivots are ``too far to the right'' given $k$.
Also marked as ``probably wrong'' in the source.

\subsection{Hash-based deduplication (\texttt{remove\_duplicates\_via\_hash})}

The RREF of a matrix (optionally with columns scaled to a canonical form) is
hashed. Matrices whose canonical RREF has been seen before are discarded. This
is disabled by default with a warning that it can exhaust memory.

\subsection{Equivalent-places symmetry reduction}

As described in Section \ref{sec:vertex-matches}, the enumeration of rows
respects the current equivalence classes of bat positions. This is not a
pruning step per se but a reorganisation of the enumeration that avoids
generating permutations of rows that are known to be equivalent.

A further (commented-out) idea in the source discusses pruning column-sign
equivalences: if two L-matrices differ only in the signs of certain columns,
and one precedes the other in the DFS, the later one need not be explored. This
is noted as potentially very powerful but ``frighteningly'' hard to justify
correctness for.

%% -----------------------------------------------------------------------
\section{Vertex-match signature enumeration}
\label{sec:signatures}

Rather than enumerating vertex matches directly, the code first enumerates
their \emph{signatures}: the triple $(e, o, z)$ giving the count of $+1$s,
$-1$s and $0$s respectively (with $e + o + z = M$, $e$ even, $o$ odd in
the tristate case). Signatures are generated in an order such that when
converted to canonical tuples (sorted non-decreasingly) the result is in
ascending lexicographic order.

For each signature, all distinct permutations of the corresponding canonical
tuple are generated, subject to the current equivalent-places constraints. This
two-level generation (signature, then permutations) is more efficient than
iterating over all $3^M$ possible tuples and filtering.

%% -----------------------------------------------------------------------
\section{Data flow summary}

The following diagram gives the overall data flow.

\medskip
\noindent
\begin{tabular}{l p{9cm}}
  \textbf{Input} & $M$, $k$, and (optionally) specific unscaled bat matrices $W$. \\[4pt]
  \textbf{Step 1} & Generate unscaled bat matrix $W$ in general position (random
    integers, sampled from a Gaussian, with a given seed). \\[4pt]
  \textbf{Step 2} & Enumerate L-matrices via DFS over vertex-match rows
    (\texttt{generate\_viable\_vertex\_match\_matrices}). \\[4pt]
  \textbf{Step 3} & For each L-matrix, form the alpha-attacking matrix $A =
    A(L, W)$ and compute $\mathcal{N}(A)$. \\[4pt]
  \textbf{Step 4} & If non-collapsing, find $\alpha^* \in \mathcal{N}(A)$ with
    all components non-zero (\texttt{nonzero\_lin\_comb}). \\[4pt]
  \textbf{Step 5} & Compute $B = \operatorname{diag}(\alpha^*) W$ and evaluate
    $(EE, OO)$ via meet-in-the-middle. \\[4pt]
  \textbf{Step 6} & Record $|EE|$ if it is the new minimum. \\[4pt]
  \textbf{Output} & The minimum observed $|EE|$, together with the L-matrix,
    scalings, and scaled bat matrix that achieved it.
\end{tabular}

%% -----------------------------------------------------------------------
\section{Notes and caveats}

\begin{itemize}
  \item All arithmetic is exact, using SymPy rational arithmetic throughout.
    This avoids floating-point errors in the null-space computation at the cost
    of speed.

  \item The output is an upper bound on the true minimum confusable-pair size
    for two distinct reasons, discussed in Section~\ref{sec:background}: (1)
    some pruning strategies may be over-aggressive, missing pairs within the
    hypercube family; and (2), more fundamentally, the hypercube construction
    is not known to cover all possible confusable pairs---the true minimum
    might arise from a completely different structure.

  \item The \texttt{only\_output\_odd\_ones} flag in \texttt{config.py}
    restricts the search to vertex matches with an odd number of $+1$ entries.
    The rationale for this restriction is not documented in the code, but may
    be related to a symmetry argument. \textbf{[?]}

  \item The \texttt{MinimalConfusableSets} subdirectory is not yet integrated
    with the rest of the Echinocoder project. The interface between it and
    ``Tom's'' decider functions (referred to in the README and
    \texttt{deciders.py}) was under design at the time the code was last
    actively developed.
\end{itemize}

\end{document}
